%% titre: Quiz 2 — Ensembles et relations
%% niveau: Licence 1
%% matiere: Algèbre 1
%% annee: 2026-2027
%% ordre: 2
%% description: Opérations sur les ensembles, produit cartésien, relations d'équivalence et d'ordre.

\begin{question}
Pour trois ensembles $A,B,C$, l'ensemble $A\cap(B\cup C)$ est égal à :
\begin{reponses}
  \bonne $(A\cap B)\cup(A\cap C)$
  \mauvaise $(A\cap B)\cap(A\cap C)$
  \mauvaise $A\cup(B\cap C)$
  \mauvaise $(A\cup B)\cap C$
\end{reponses}
\begin{explication}
C'est la distributivité de $\cap$ sur $\cup$ : $x\in A$ et ($x\in B$ ou $x\in C$).
\end{explication}
\end{question}

\begin{question}
Dans un ensemble $E$, le complémentaire de $A\cup B$ est :
\begin{reponses}
  \bonne $A^{c}\cap B^{c}$
  \mauvaise $A^{c}\cup B^{c}$
  \mauvaise $A^{c}\setminus B^{c}$
  \mauvaise $A\cap B$
\end{reponses}
\begin{explication}
Loi de De Morgan : $(A\cup B)^c=A^c\cap B^c$.
\end{explication}
\end{question}

\begin{question}
Si $E$ est un ensemble fini à $n$ éléments, l'ensemble $\mathcal{P}(E)$ de ses parties a pour cardinal :
\begin{reponses}
  \mauvaise $n^2$
  \mauvaise $2n$
  \bonne $2^n$
  \mauvaise $n!$
\end{reponses}
\begin{explication}
Chaque élément est soit dans la partie, soit en dehors : $2$ choix par élément, donc $2^n$ parties.
\end{explication}
\end{question}

\begin{question}
Quelles propriétés possède la relation $\le$ sur $\R$ ? (plusieurs réponses possibles)
\begin{reponses}
  \bonne réflexive
  \bonne antisymétrique
  \bonne transitive
  \mauvaise symétrique
\end{reponses}
\begin{explication}
Réflexive, antisymétrique et transitive : c'est une relation d'ordre. Elle n'est pas symétrique ($1\le2$ mais $2\not\le1$).
\end{explication}
\end{question}

\begin{question}
Sur $\Z$, on pose $x\,\mathcal{R}\,y \iff x-y\in 3\Z$. Combien cette relation d'équivalence a-t-elle de classes ?
\begin{reponses}
  \mauvaise $1$
  \mauvaise $2$
  \bonne $3$
  \mauvaise une infinité
\end{reponses}
\begin{explication}
Deux entiers sont en relation s'ils ont le même reste dans la division par $3$ : les classes sont celles des restes $0$, $1$ et $2$.
\end{explication}
\end{question}

\begin{question}
Quelles familles sont des partitions de $E=\{1,2,3,4\}$ ? (plusieurs réponses possibles)
\begin{reponses}
  \bonne $\big\{\{1,2\},\{3,4\}\big\}$
  \bonne $\big\{\{1\},\{2,3,4\}\big\}$
  \mauvaise $\big\{\{1,2\},\{2,3\},\{4\}\big\}$
  \mauvaise $\big\{\{1\},\{2\},\{3\}\big\}$
  \mauvaise $\big\{\varnothing,\{1,2,3,4\}\big\}$
\end{reponses}
\begin{explication}
Une partition est formée de parties non vides, deux à deux disjointes, dont la réunion est $E$. La troisième famille a des parties qui se chevauchent, la quatrième oublie $4$, la dernière contient $\varnothing$.
\end{explication}
\end{question}

\begin{question}
Si $A$ et $B$ sont finis, le cardinal de $A\times B$ vaut :
\begin{reponses}
  \mauvaise $\operatorname{Card}A+\operatorname{Card}B$
  \bonne $\operatorname{Card}A\times\operatorname{Card}B$
  \mauvaise $\max(\operatorname{Card}A,\operatorname{Card}B)$
  \mauvaise $\operatorname{Card}A^{\operatorname{Card}B}$
\end{reponses}
\end{question}
